\chapter{Описание компании и исходных данных}

\section{Характеристика компании}

Исследуемая компания занимается производством пищевой продукции: сладостей, печенья и выпечки. Основные каналы продаж -- крупные сети дискаунтеров, магазины низких цен и торговые рынки. У предприятия есть два производственных участка, которые связаны между собой и дополняют друг друга.

Ключевая управленческая проблема заключается в том, что бизнес функционирует, оплачивает зарплаты, кредиты и обязательства, но не формирует устойчивый свободный денежный поток для вывода средств собственниками.

\section{Источники данных}

В проекте используются следующие группы данных:
\begin{itemize}
  \item банковские выписки по расчетному счету;
  \item продажи по контрагентам и номенклатуре;
  \item продажи по оплатам;
  \item корректировки реализации и возвраты;
  \item справочник товаров;
  \item результаты ручной проверки операций.
\end{itemize}

Сырые данные хранятся в виде Excel-отчетов и банковских выгрузок. Для анализа они приводятся к итоговым таблицам: \texttt{sales\_by\_counterparties}, \texttt{sales\_by\_payment\_counterparties} и \texttt{bank\_statement\_transactions}.

\section{Пайплайн подготовки данных}

Подготовка данных построена как воспроизводимый процесс. Сырые файлы помещаются в соответствующие папки, затем запускаются Python-пайплайны, которые очищают данные, унифицируют поля и сохраняют результат в форматах Parquet, CSV и XLSX.

Для каждого источника сохраняются служебные поля \texttt{source\_*}, которые позволяют отследить происхождение строки: исходный файл, лист, номер строки и хэш файла. Это важно для аудита данных и повторной обработки при изменении исходных выгрузок.

\section{Первичный анализ качества данных}

Перед применением ML необходимо проверить техническое качество данных: корректность дат, наличие дублей, пропуски, пересекающиеся выписки, разные форматы сумм и неоднородные названия контрагентов. В проекте отдельное внимание уделено исправлению парсинга дат и удалению только технических дублей, чтобы не потерять реальные операции.
